\documentclass[14 pt]{extarticle}

	\usepackage[frenchb]{babel}
	\usepackage[utf8]{inputenc}  
	\usepackage[T1]{fontenc}
	\usepackage{amssymb}
	\usepackage[mathscr]{euscript}
	\usepackage{stmaryrd}
	\usepackage{amsmath}
	\usepackage{tikz}
	\usepackage[all,cmtip]{xy}
	\usepackage{amsthm}
	\usepackage{varioref}
	\usepackage{geometry}
	\geometry{a4paper}
	\usepackage{lmodern}
	\usepackage{hyperref}
	\usepackage{array}
	 \usepackage{fancyhdr}
\renewcommand{\theenumi}{\alph{enumi})}
	\pagestyle{fancy}
	\theoremstyle{plain}
	\fancyfoot[C]{} 
	\fancyhead[L]{Contrôle}
	\fancyhead[R]{octobre 2025}\geometry{
 a4paper,
 total={170mm,257mm},
 left=20mm,
 top=20mm,
 }
	
	
	\title{Chapitre 1-  Arithmétique}
	\date{}
	\begin{document}

\begin{center}{\Large Contrôle Chapitre 1 - Arithmétique}\\ 
 \end{center} 
 
 \textbf{Toute réponse doit être dûment détaillée. La calculcatrice est interdite.}
\subsection*{Exercice 1 - 6 points} 
 
 \begin{enumerate}
 \item Poser la division euclidienne de 12~738 par 6. 
 \item Le nombre $441$ est-il premier ? 
 \item Trouver tous les diviseurs de $90$. 
 \end{enumerate}
 
\subsection*{Exercice 2 - 5 points} 
\begin{enumerate}
\item Décomposer $84$ et $360$ en facteurs premiers. 
\item Faire la liste des diviseurs de $84$. 
\item Trouver le plus grand diviseur commun de $84$ et $360$.
\end{enumerate} 

 
 \subsection*{Exercice 3 (Métropole, 2024) - 5 points}
 
La présidente du club veut offrir des petits sachets cadeaux tous identiques contenant des autocollants et des drapeaux avec le logo du club. Elle a acheté 330 autocollants et 132 drapeaux et veut tous les utiliser. Elle veut que, dans chaque sachet, il y ait exactement le même nombre d'autocollants et que, dans chaque sachet, il y ait exactement le même nombre de drapeaux.

\medskip

\begin{enumerate}
\item Pourquoi n'est-il pas possible de faire 15 sachets ?
\item 
	\begin{enumerate}
		\item Décomposer 330 et 132 en produits de facteurs premiers.
		\item En déduire le plus grand nombre de sachets que la présidente pourra réaliser.
		\item Dans ce cas, combien mettra-t-elle d'autocollants et de drapeaux dans chaque sachet ?
	\end{enumerate}
\end{enumerate}
 
 \subsection*{Exercice 4 - 3 points}
 \begin{enumerate}
 \item 
 Simplifier au maximum la fraction $\frac{420}{1008}$.
 \item 
 Même question avec la fraction $\frac{345}{2205}$.
 \end{enumerate}
 \subsection*{Exercice 5 - 1 point}
  Recopier la bonne décomposition en facteurs premiers de 987~654~321 : 
  \[ 3^2\times 17^2\times 379721 ;\phantom{meow} 2^3 \times 17 \times 379721 ;\phantom{meow} 2^3 \times 17 \times 379722 \]
 \newpage\begin{center}{\Large Contrôle Chapitre 1 - Arithmétique}\\ 
 \end{center} 
 
 \textbf{Toute réponse doit être dûment détaillée. La calculcatrice est interdite.}
\subsection*{Exercice 1 - 6 points} 
 
 \begin{enumerate}
 \item Poser la division euclidienne de 24~438 par 6. 
 \item Le nombre $431$ est-il premier ? 
 \item Trouver tous les diviseurs de $150$. 
 \end{enumerate}
 
\subsection*{Exercice 2 - 5 points} 
\begin{enumerate}
\item Décomposer $252$ et $342$ en facteurs premiers. 
\item Faire la liste des diviseurs de $342$. 
\item Trouver le plus grand diviseur commun de $252$ et $342$.
\end{enumerate} 

 
 \subsection*{Exercice 3 (Métropole, 2024) - 5 points}
 
La présidente du club veut offrir des petits sachets cadeaux tous identiques contenant des autocollants et des drapeaux avec le logo du club. Elle a acheté 330 autocollants et 132 drapeaux et veut tous les utiliser. Elle veut que, dans chaque sachet, il y ait exactement le même nombre d'autocollants et que, dans chaque sachet, il y ait exactement le même nombre de drapeaux.

\medskip

\begin{enumerate}
\item Pourquoi n'est-il pas possible de faire 15 sachets ?
\item 
	\begin{enumerate}
		\item Décomposer 330 et 132 en produits de facteurs premiers.
		\item En déduire le plus grand nombre de sachets que la présidente pourra réaliser.
		\item Dans ce cas, combien mettra-t-elle d'autocollants et de drapeaux dans chaque sachet ?
	\end{enumerate}
\end{enumerate}
 
 \subsection*{Exercice 4 - 3 points}
 \begin{enumerate}
 \item 
 Simplifier au maximum la fraction $\frac{294}{462}$.
 \item 
 Même question avec la fraction $\frac{345}{2205}$.
 \end{enumerate}
 \subsection*{Exercice 5 - 1 point}
  Recopier la bonne décomposition en facteurs premiers de 987~654~321 : 
  \[ 3\times5\times 3~607\times 3~803 ;\phantom{meow} 3^2\times 3~607\times 3~803 ;\phantom{meow} 3\times 10~821\times 3~803 \]
 
 	\end{document}
